\documentclass[11pt,letterpaper]{article}
\def\SuiteDocTitle{Acceptance Criteria and Test Engineering}
\def\SuiteDocCode{TR-03}
% Shared professional style: pdfLaTeX compatible, no external font files.
\usepackage[T1]{fontenc}
\usepackage[utf8]{inputenc}
\usepackage{lmodern}
\usepackage{microtype}
\usepackage[margin=0.83in,headheight=16pt,footskip=28pt]{geometry}
\usepackage[table]{xcolor}
\definecolor{navy}{HTML}{12304A}
\definecolor{teal}{HTML}{087F83}
\definecolor{gold}{HTML}{B78633}
\definecolor{blueback}{HTML}{F0F5F8}
\definecolor{tealback}{HTML}{EEF8F6}
\definecolor{amberback}{HTML}{FFF7E8}
\definecolor{grayback}{HTML}{F6F7F8}
\usepackage{array,tabularx,longtable,booktabs}
\usepackage{enumitem}
\usepackage{amssymb}
\usepackage{titlesec}
\usepackage{fancyhdr,lastpage}
\usepackage{tcolorbox}
\tcbuselibrary{breakable}
\usepackage{pgffor}
\usepackage{hyperref}
\hypersetup{colorlinks=true,linkcolor=teal,urlcolor=teal,pdfauthor={Engineering Learning Series},pdftitle={\SuiteDocTitle}}
\setlength{\parindent}{0pt}
\setlength{\parskip}{5.5pt}
\setlist{topsep=3pt}
\setcounter{tocdepth}{2}
\titleformat{\section}{\Large\bfseries\color{navy}}{\thesection}{0.7em}{}
\titleformat{\subsection}{\large\bfseries\color{teal}}{\thesubsection}{0.6em}{}
\titleformat{\subsubsection}{\normalsize\bfseries\color{navy}}{\thesubsubsection}{0.6em}{}
\pagestyle{fancy}
\fancyhf{}
\fancyhead[L]{\scriptsize\color{navy}AGILE REQUIREMENTS ENGINEERING}
\fancyhead[R]{\scriptsize\color{teal}\SuiteDocCode}
\fancyfoot[L]{\scriptsize Independent learning and implementation materials}
\fancyfoot[R]{\scriptsize Page \thepage\ of \pageref{LastPage}}
\renewcommand{\headrulewidth}{0.35pt}
\renewcommand{\footrulewidth}{0pt}
\newcommand{\SuiteHeader}[2]{%
{\color{teal}\small\bfseries ENGINEERING DOCUMENTATION SUITE \hfill #2}\par\vspace{8pt}
{\color{navy}\LARGE\bfseries\SuiteDocTitle}\par\vspace{5pt}
{\color{black!68}\normalsize #1}\par\vspace{8pt}
{\color{gold}\rule{\linewidth}{2pt}}\par\vspace{8pt}}
\newcommand{\WritingLines}[1]{\par\foreach \n in {1,...,#1}{\noindent\rule{\linewidth}{0.35pt}\par\vspace{0.28cm}}}
\newcommand{\FreeSpace}[1]{\begin{tcolorbox}[colback=white,colframe=black!25,height=#1,boxrule=0.5pt]\end{tcolorbox}}
\newcommand{\Field}[1]{\noindent\textbf{#1}\quad\rule{0.42\linewidth}{0.35pt}\par\vspace{0.3cm}}
\newcommand{\CheckItem}[1]{\noindent$\square$\quad #1\par\vspace{2pt}}
\newcommand{\SmallNote}[1]{{\footnotesize\color{black!62}#1}\par}
\newcommand{\Guidance}[1]{\begin{tcolorbox}[colback=blueback,colframe=teal,breakable,title=Guidance]#1\end{tcolorbox}}
\sloppy

\begin{document}
\SuiteHeader{Specification by example | executable evidence and test ownership}{Reference}
\tableofcontents\medskip\noindent\rule{\linewidth}{0.4pt}\medskip
\section{Principles and roles}
Behavioral acceptance criteria clarify what the product needs; acceptance tests produce evidence that the system behaves accordingly. Product representatives determine desired outcomes, testers make examples discriminating, and engineers automate and investigate failure. Create the examples before the implementation when practical, then adjust them through shared conversation.

\section{From intent to executable scenarios}
\begin{enumerate}[leftmargin=*,itemsep=3pt]
\item Identify the actor, trigger, permission and expected business result.
\item Choose a concrete starting state and explicit sample data.
\item Specify one action and at least one externally observable outcome.
\item Add exceptional paths, denied actions and state transitions that change the business meaning.
\item Link examples to an implementation and verification layer, then record evidence.
\end{enumerate}
\section{Illustrative BDD-style examples}
\begin{tcolorbox}[colback=grayback,colframe=navy,title={Original example scenarios},breakable]
\ttfamily\small
Feature: Service request submission\\
Scenario: Valid request\\
Given an authenticated requester with a valid request category\\
And a complete request description\\
When the requester submits the request\\
Then the system records a new request\\
And displays a unique tracking reference\\[4pt]
Scenario: Denied cross-account access\\
Given a request owned by a different requester\\
When the requester attempts to open that request\\
Then the system denies access\\
And reveals no request details
\end{tcolorbox}
\section{Test-design matrix}
\begin{longtable}{|p{0.20\linewidth}|p{0.29\linewidth}|p{0.37\linewidth}|}\hline
\textbf{Class} & \textbf{Question} & \textbf{Sample} \\ \hline\endfirsthead\hline
\textbf{Class} & \textbf{Question} & \textbf{Sample} \\ \hline\endhead
Positive & Can an eligible actor accomplish the goal? & Valid request is accepted \\ \hline
Negative & Is an invalid action rejected safely? & Missing mandatory reason is rejected \\ \hline
Boundary & What happens at permitted limits? & Maximum-length input is handled \\ \hline
Authorization & Can an actor see or change another actor’s records? & Cross-tenant read is denied \\ \hline
Transition & Are state changes legal and persistent? & Approved item cannot be silently reopened \\ \hline
Quality & Does a measurable property hold? & Latency target under test profile \\ \hline
Failure and recovery & Is failure observable and reversible? & Retry does not create duplicate request \\ \hline
\end{longtable}
\section{Automation and evidence}
Use stable scenario IDs, version-controlled test definitions, deterministic test fixtures, isolated environments and clean teardown. Evidence should identify build revision, scenario, execution time, result, relevant log reference and defect link. Avoid leaking real personal data into fixtures. Test design must also include human-in-the-loop usability and accessibility validation where automation cannot establish suitability.

\section{Failure triage}
\begin{enumerate}[leftmargin=*,itemsep=3pt]
\item Determine whether the failure is product behavior, flaky environment or invalid test expectation.
\item Confirm the expected behavior with the product representative when it is ambiguous.
\item Record reproducible inputs and environment, then fix code or renegotiate requirement.
\item Keep the failed test or equivalent regression protection when the behavior remains required.
\item Attach final evidence to story acceptance and release review.
\end{enumerate}
\section{Source boundary}
Cohn, Chapter 6 and Chapter 21, provide the book-grounded emphasis on writing acceptance tests early and using them to convey customer expectations. Gherkin-style scenarios, CI evidence metadata and security automation structures here are independently authored extensions.


\end{document}
